\chapter{\abstractname}

Standard multilingual evaluation benchmarks often conflate internal representation quality with prompt sensitivity, reasoning capacity, and generation dynamics. Under the pivot hypothesis, multilingual transformers develop a shared, language-agnostic representation space in their intermediate layers. This thesis examines internal cross-lingual alignment using the MEXA framework, a reference-free method that extracts layer-wise representations, applies position-weighted pooling, and evaluates bidirectional Precision-at-1 retrieval on parallel sentences.

We test 16 architectures across dense causal decoders (Qwen~3, Qwen~3.5~9B, Apertus, Llama~3.1~8B, Mistral~7B), sparse Mixture-of-Experts (Mixtral~8x7B and 8x22B), dedicated sentence encoders (LaBSE, mE5), and masked language models across FLORES-200 and the 1{,}401-language Super-Parallel Bible Corpus. Parameter scaling and diverse pre-training data raise peak alignment up to $\mu_{\text{Max}} = 0.7794$ (Qwen~3.5~9B), and language difficulty rankings remain consistent across domains ($\rho = 0.98$). In sparse MoE models, routing depends on syntactic and semantic roles rather than language identity, keeping the shared interlingua intact. Evaluating non-English pivots (French, German, Arabic, Chinese, Basque) shows that English is not an obligatory hub; mean-centering removes language centroid offsets and brings peak alignment across all pivots to $\mu_{\text{Max}} \ge 0.94$. Although high tokenizer fertility penalizes low-resource scripts ($r = -0.56$), MEXA scores predict zero-shot task accuracy on Belebele and m-ARC ($\rho = 1.00$). Together, these findings show that training-free representation alignment serves as an accurate, low-cost diagnostic for multilingual foundation models.
